% ============================================
% Chapter 7 – Conclusion and Future Work
% ============================================

\chapter{Conclusion and Future Work}
\justifying

% -------------------------------------------------
% 7.1 Conclusion
% -------------------------------------------------

\section{Conclusion}

Overall, the "Lets Go Ride Sharing App" successfully tackled the major problems of urban transportation in Pakistan, such as cost-effectiveness, security, and trust that can be built around the platform through community participation. The system's many-to-many ride sharing function connects private car owners and everyday commuters and makes a viable alternative to the conventional ride-hailing service.

In terms of technology, the project's primary goals were met with a solid front-end framework, Flutter, for cross-platform apps and a powerful and secure backend, Django.From a technical standpoint, it successfully integrated a robust front-end framework, Flutter, for cross-platform application development, along with a powerful and secure back-end framework, Django. A content-based recommendation engine helps users connect with the rides they've historically enjoyed and/or are gender comfortable with, and geospatial geofencing for OTP verification adds a concrete level of security when executing an OTP.

The reliability of the system was verified during the evaluation stage, having 480 automated test cases, of which 281 were in the frontend and 199 in the backend, passed it successfully with a pass rate of 100\%. User acceptance testing also confirmed that the interface is user-friendly and the social elements, which include clear negotiation and SOS networking, are very appreciated by the target community.

Overall, "Lets Go Ride Sharing App" is a promising and scalable idea that harnesses current AI and mobile technologies to foster a more connected and secure commuting experience.

% -------------------------------------------------
% 7.2 Future Work
% -------------------------------------------------

\section{Future Work}

The current iterations of "Lets Go Ride Sharing App" serves as a strong structure for community ride sharing, though there are a number of avenues to expand upon:

\begin{itemize}
    \item \textbf{Integrated Payment Gateways:} It will move toward cashless payments with integrated payment gateways like EasyPaisa, JazzCash, and Stripe in the future.
    \item To further improve social safety, a filter could be set up that allows for "Ladies-Only" trips, giving women the opportunity to choose to take a ride car to car with other women.
    \item By incorporating real-time traffic data from APIs like OSRM or Google Maps into the custom-built ranking mechanism, the algorithm could dynamically update ETA based on traffic conditions.
    \item The introduction of a module for calculating CO$_2$ emissions avoided by the use of ride-sharing would raise awareness of environmental issues and provide "Green Points", rewards to the users.
    \item Enabling corporate fleet use would enable businesses to optimize their employee commute management.
\end{itemize}

% -------------------------------------------------
% 7.3 Limitations
% -------------------------------------------------

\section{Limitations}

However, the project currently has the following restrictions:

\begin{itemize}
    \item \textbf{Data Connectivity:} It is assumed that the system is connected to a stable internet connection at all times with the possibility of receiving updates every hour or so via HTTP. No offline map support is implemented at this time.
    \item \textbf{Manual Verification Overhead:} Manual check on driver documents (CNIC, License) in Admin Dashboard is a bottleneck for scalability which needs to be automated by OCR in future phase.
    \item \textbf{The Cold-Start Problem:} First, the new users are given suggestions based on the popularity.First, the new users are given suggestions based on the popularity (Cold-Start Problem).
\end{itemize}
